\chapter{Về Suy Nghĩ}
\markboth{Về Suy Nghĩ}{About Thinking}

\section{Đừng Nghĩ Mình Ngu}
\begin{truyen}{Đừng Nghĩ Mình Ngu}{Do Not Think You Are Stupid}
\chuhoa{N}{hiều học sinh chỉ vì vài lần chưa hiểu bài đã vội kết luận mình ngu.}
\textit{(Many students quickly conclude they are stupid after only a few moments of confusion.)}

Cách nghĩ đó rất nặng và thường không đúng.
\textit{(That thought is heavy and often untrue.)}

Nó làm em mất động lực trước cả khi thật sự cố gắng đủ lâu.
\textit{(It steals motivation before you have truly tried for long enough.)}

Thay vì nghĩ mình ngu, em nên nghĩ rằng mình còn đang học và còn cần cách phù hợp hơn.
\textit{(Instead of calling yourself stupid, think of yourself as still learning and still needing a better way.)}
\end{truyen}

\begin{baihoc}
Đừng biến một giai đoạn khó thành bản án cho trí óc của mình.
\textit{(Do not turn a difficult period into a sentence on your mind.)}
\end{baihoc}

\begin{ghinhoanh}
\textbf{Short English to remember:}

Temporary difficulty is not permanent stupidity.

\end{ghinhoanh}

\begin{tuhocnhanh}
1. Vì sao không nên nghĩ mình ngu? \textit{Why should you not call yourself stupid?}

2. Cách nghĩ đó hại điều gì? \textit{What does that way of thinking damage?}

3. Em có thể thay nó bằng câu nào công bằng hơn? \textit{What fairer sentence can replace it?}
\end{tuhocnhanh}
\ngancach

\section{Đừng Nghĩ Mình Giỏi Quá}
\begin{truyen}{Đừng Nghĩ Mình Giỏi Quá}{Do Not Think You Are Too Good Already}
\chuhoa{T}{ự tin là cần, nhưng tự mãn rất nguy hiểm.}
\textit{(Confidence is necessary, but arrogance is dangerous.)}

Khi nghĩ mình giỏi quá rồi, học sinh thường ít nghe góp ý và ít chịu sửa.
\textit{(When students think they are already very good, they listen less and correct less.)}

Điều đó làm sự tiến bộ chậm lại trong khi bản thân vẫn tưởng mình ổn.
\textit{(That slows progress while the student still feels fine.)}

Người học tốt lâu dài thường giữ được cả tự tin lẫn khiêm tốn.
\textit{(Strong long-term learners often keep both confidence and humility.)}
\end{truyen}

\begin{baihoc}
Tự tin giúp em đi lên; tự mãn làm em đứng lại.
\textit{(Confidence helps you move upward; arrogance keeps you still.)}
\end{baihoc}

\begin{ghinhoanh}
\textbf{Short English to remember:}

Confidence grows best with humility.

\end{ghinhoanh}

\begin{tuhocnhanh}
1. Vì sao không nên nghĩ mình giỏi quá? \textit{Why should you not think you are already too good?}

2. Tự mãn làm chậm điều gì? \textit{What does arrogance slow down?}

3. Em đang tự tin hay tự mãn? \textit{Are you confident or becoming arrogant?}
\end{tuhocnhanh}
\ngancach

\section{Suy Nghĩ Tích Cực}
\begin{truyen}{Suy Nghĩ Tích Cực}{Think Positively}
\chuhoa{S}{uy nghĩ tích cực không phải là giả vờ mọi thứ đều ổn.}
\textit{(Positive thinking is not pretending everything is fine.)}

Nó là nhìn khó khăn mà vẫn tin có cách để đi tiếp.
\textit{(It means seeing difficulty while still believing there is a way forward.)}

Một đầu óc chỉ quen nghĩ tiêu cực sẽ rất nhanh mệt và nhanh bỏ.
\textit{(A mind trained only in negativity becomes tired and gives up quickly.)}

Tích cực đúng cách giúp em bền hơn trước những đoạn chưa dễ.
\textit{(Healthy positivity helps you remain steady through difficult stretches.)}
\end{truyen}

\begin{baihoc}
Tích cực không xóa khó khăn, nhưng nó cho mình cách đứng trước khó khăn tốt hơn.
\textit{(Positive thinking does not erase difficulty, but it changes how we stand before it.)}
\end{baihoc}

\begin{ghinhoanh}
\textbf{Short English to remember:}

Healthy positivity supports endurance.

\end{ghinhoanh}

\begin{tuhocnhanh}
1. Suy nghĩ tích cực thật sự là gì? \textit{What is healthy positive thinking?}

2. Suy nghĩ tiêu cực kéo mình đi đâu? \textit{Where does negative thinking pull us?}

3. Em có thể đổi góc nhìn ở việc nào? \textit{Where can you change your perspective?}
\end{tuhocnhanh}
\ngancach

\section{Không Sợ Sai}
\begin{truyen}{Không Sợ Sai}{Do Not Fear Mistakes}
\chuhoa{N}{ỗi sợ sai làm nhiều học sinh im lặng và đứng yên.}
\textit{(Fear of mistakes makes many students silent and still.)}

Trong khi đó, lớp học là nơi con người được quyền sai để rồi sửa.
\textit{(Yet the classroom is where people are allowed to be wrong and then correct themselves.)}

Nếu sợ sai quá mức, em sẽ bỏ mất cơ hội thử và hiểu sâu hơn.
\textit{(If you fear mistakes too much, you lose the chance to try and understand more deeply.)}

Không sợ sai không có nghĩa là thích sai; nó là dám học qua sai.
\textit{(Not fearing mistakes does not mean liking them; it means daring to learn through them.)}
\end{truyen}

\begin{baihoc}
Học thật cần một chút can đảm để chấp nhận mình chưa đúng ngay.
\textit{(Real learning needs some courage to accept not being right yet.)}
\end{baihoc}

\begin{ghinhoanh}
\textbf{Short English to remember:}

Do not let fear block learning.

\end{ghinhoanh}

\begin{tuhocnhanh}
1. Vì sao không nên sợ sai quá mức? \textit{Why should you not fear mistakes too much?}

2. Không sợ sai nghĩa là gì? \textit{What does not fearing mistakes mean?}

3. Em đang ngại thử điều gì vì sợ sai? \textit{What are you avoiding because of fear?}
\end{tuhocnhanh}
\ngancach

\section{Không Sợ Thử}
\begin{truyen}{Không Sợ Thử}{Do Not Fear Trying}
\chuhoa{N}{hiều điều trong học tập chỉ mở ra khi em chịu làm thử.}
\textit{(Many things in learning open only when you are willing to try.)}

Nếu cứ đợi chắc chắn rồi mới bắt đầu, em sẽ bỏ lỡ rất nhiều cơ hội tập.
\textit{(If you wait for certainty before starting, you will miss many chances to practice.)}

Thử không bảo đảm thành công ngay.
\textit{(Trying does not guarantee immediate success.)}

Nhưng không thử gần như bảo đảm rằng em sẽ đứng yên.
\textit{(But not trying almost guarantees that you will stay still.)}
\end{truyen}

\begin{baihoc}
Nhiều cánh cửa của việc học không mở cho người chỉ đứng nhìn.
\textit{(Many doors of learning do not open for those who only stand and look.)}
\end{baihoc}

\begin{ghinhoanh}
\textbf{Short English to remember:}

Trying creates possibility.

\end{ghinhoanh}

\begin{tuhocnhanh}
1. Vì sao cần không sợ thử? \textit{Why is it important not to fear trying?}

2. Đợi chắc chắn quá lâu gây hại gì? \textit{What harm comes from waiting too long for certainty?}

3. Em cần thử điều gì trước? \textit{What do you need to try first?}
\end{tuhocnhanh}
\ngancach

\section{Không Sợ Hỏi}
\begin{truyen}{Không Sợ Hỏi}{Do Not Fear Asking}
\chuhoa{H}{ỏi là một việc rất bình thường của người đang học.}
\textit{(Asking is a very normal action for someone who is learning.)}

Ngại hỏi lâu ngày khiến nhiều chỗ không hiểu trở thành nỗi sợ lớn hơn.
\textit{(Avoiding questions for too long turns confusion into bigger fear.)}

Không sợ hỏi giúp em giữ việc học gần với sự thật hơn.
\textit{(Not fearing questions keeps learning closer to reality.)}

Nó cũng cho thấy em coi việc hiểu thật quan trọng hơn vẻ ngoài biết rồi.
\textit{(It shows that you value real understanding more than appearances.)}
\end{truyen}

\begin{baihoc}
Một câu hỏi đúng lúc có thể cứu rất nhiều thời gian và nhiều nỗi ngại về sau.
\textit{(One timely question can save much time and later embarrassment.)}
\end{baihoc}

\begin{ghinhoanh}
\textbf{Short English to remember:}

Questions support honest learning.

\end{ghinhoanh}

\begin{tuhocnhanh}
1. Vì sao không nên sợ hỏi? \textit{Why should you not fear asking?}

2. Hỏi giúp giữ việc học gần với điều gì? \textit{What does asking keep learning close to?}

3. Em có câu hỏi nào nên nói ra không? \textit{Is there a question you should say out loud?}
\end{tuhocnhanh}
\ngancach

\section{Không Sợ Bắt Đầu}
\begin{truyen}{Không Sợ Bắt Đầu}{Do Not Fear Beginning}
\chuhoa{N}{hiều học sinh không hẳn lười, mà chỉ sợ bắt đầu một việc còn mờ và khó.}
\textit{(Many students are not truly lazy; they are simply afraid to begin something unclear and difficult.)}

Nhưng mọi hành trình đều cần một bước đầu, dù còn vụng và nhỏ.
\textit{(But every journey needs a first step, even if small and awkward.)}

Khi đã bắt đầu, đầu óc thường dễ đi tiếp hơn ta tưởng.
\textit{(Once begun, the mind often continues more easily than expected.)}

Không sợ bắt đầu là một sức mạnh rất thiết thực.
\textit{(Not fearing the beginning is a very practical strength.)}
\end{truyen}

\begin{baihoc}
Điều khó nhất nhiều khi không phải làm tiếp, mà là dám đi vào bước đầu tiên.
\textit{(Often the hardest part is not continuing, but daring to take the first step.)}
\end{baihoc}

\begin{ghinhoanh}
\textbf{Short English to remember:}

Starting often breaks more fear than waiting.

\end{ghinhoanh}

\begin{tuhocnhanh}
1. Vì sao nhiều học sinh sợ bắt đầu? \textit{Why do many students fear beginning?}

2. Bắt đầu rồi thường có điều gì thay đổi? \textit{What often changes after you begin?}

3. Em đang trì hoãn bắt đầu việc nào? \textit{What are you postponing starting?}
\end{tuhocnhanh}
\ngancach

\section{Tin Vào Nỗ Lực}
\begin{truyen}{Tin Vào Nỗ Lực}{Believe in Effort}
\chuhoa{N}{ỗ lực không bảo đảm mọi điều sẽ đến rất nhanh, nhưng nó là điều giữ cho khả năng còn sống.}
\textit{(Effort does not guarantee quick success, but it keeps possibility alive.)}

Nếu không tin vào nỗ lực, học sinh sẽ dễ buông xuôi trước cả khi sự thay đổi kịp đến.
\textit{(If you do not believe in effort, you give up before change has time to arrive.)}

Tin vào nỗ lực không phải là mù quáng.
\textit{(Believing in effort is not blind optimism.)}

Nó là chấp nhận rằng có những điều chỉ cải thiện khi mình chịu làm phần của mình.
\textit{(It means accepting that some things improve only when we do our part.)}
\end{truyen}

\begin{baihoc}
Nỗ lực có thể chậm, nhưng thiếu nó thì rất khó có điều gì khá lên thật.
\textit{(Effort may be slow, but without it very little truly improves.)}
\end{baihoc}

\begin{ghinhoanh}
\textbf{Short English to remember:}

Effort keeps improvement possible.

\end{ghinhoanh}

\begin{tuhocnhanh}
1. Vì sao cần tin vào nỗ lực? \textit{Why is it important to believe in effort?}

2. Tin vào nỗ lực khác gì với mù quáng? \textit{How is this different from being blind?}

3. Em đang cần kiên trì với nỗ lực nào? \textit{What effort do you need to stay with?}
\end{tuhocnhanh}
\ngancach

\section{Tự Suy Nghĩ}
\begin{truyen}{Tự Suy Nghĩ}{Think for Yourself}
\chuhoa{H}{ọc giỏi không chỉ là biết nhắc lại, mà còn là biết tự nghĩ.}
\textit{(Learning well is not only about repeating, but also about thinking for yourself.)}

Có những câu hỏi không thể trả lời tử tế nếu đầu óc chỉ chờ sẵn đáp án của người khác.
\textit{(Some questions cannot be answered well if the mind only waits for others' ready-made answers.)}

Tự suy nghĩ giúp em hiểu sâu hơn và đứng vững hơn trong cách nhìn của mình.
\textit{(Thinking for yourself helps you understand more deeply and stand more firmly in your own view.)}

Đó là một năng lực rất đáng luyện từ sớm.
\textit{(It is a skill worth training early.)}
\end{truyen}

\begin{baihoc}
Đầu óc mạnh lên khi được dùng để nghĩ thật, không chỉ để chép lại.
\textit{(The mind grows stronger when used for real thinking, not only repetition.)}
\end{baihoc}

\begin{ghinhoanh}
\textbf{Short English to remember:}

Independent thinking builds real understanding.

\end{ghinhoanh}

\begin{tuhocnhanh}
1. Vì sao cần tự suy nghĩ? \textit{Why is independent thinking important?}

2. Chờ sẵn đáp án của người khác có hại gì? \textit{What is harmful about waiting only for others' answers?}

3. Em đang tự nghĩ nhiều hay chỉ lặp lại? \textit{Are you thinking independently or mostly repeating?}
\end{tuhocnhanh}
\ngancach

\section{Luôn Tò Mò}
\begin{truyen}{Luôn Tò Mò}{Stay Curious}
\chuhoa{T}{ò mò là động cơ rất đẹp của việc học.}
\textit{(Curiosity is a beautiful engine for learning.)}

Nó khiến em muốn biết thêm, hỏi thêm và nhìn rộng hơn những gì đang có trước mắt.
\textit{(It makes you want to know more, ask more, and see farther than what is in front of you.)}

Khi mất tò mò, việc học rất dễ chỉ còn là nghĩa vụ.
\textit{(When curiosity is lost, learning easily becomes only duty.)}

Giữ được sự tò mò là giữ được một phần rất sống của trí óc.
\textit{(Keeping curiosity means keeping a very alive part of the mind.)}
\end{truyen}

\begin{baihoc}
Người còn tò mò thường còn học được rất nhiều.
\textit{(Those who remain curious can still learn a great deal.)}
\end{baihoc}

\begin{ghinhoanh}
\textbf{Short English to remember:}

Curiosity keeps the mind alive.

\end{ghinhoanh}

\begin{tuhocnhanh}
1. Vì sao nên luôn tò mò? \textit{Why should you stay curious?}

2. Khi mất tò mò, việc học dễ thành gì? \textit{What does learning become when curiosity fades?}

3. Em đang tò mò muốn hiểu thêm điều gì? \textit{What are you curious to understand more deeply?}
\end{tuhocnhanh}
\ngancach